\section*{Part A: Solved Examples 1 to 5 (Section 7.2, Distance Formula)}

%---------------------------------------------------------------
\subsection*{Example 1}
\textbf{Question.} Do the points $(3,2)$, $(-2,-3)$ and $(2,3)$ form a triangle? If so, name the type of triangle formed.

\textbf{Solution.} Let $P(3,2)$, $Q(-2,-3)$ and $R(2,3)$.

\textbf{Step 1: find the three sides.} By the distance formula, the distance between $(x_1,y_1)$ and $(x_2,y_2)$ is
\[d=\sqrt{(x_2-x_1)^2+(y_2-y_1)^2}.\]
\distnote
\begin{align*}
\dist{PQ}{3}{2}{-2}{-3}{5\sqrt2\approx 7.07}\\
\dist{QR}{-2}{-3}{2}{3}{2\sqrt{13}\approx 7.21}\\
\distn{PR}{3}{2}{2}{3}\approx 1.41
\end{align*}
\slip{A negative number squared is positive: $(-5)^2=25$, not $-25$. Also $-2-3=-5$ (not $1$), and $2-(-2)=2+2=4$.}

\textbf{Step 2: do the points form a triangle?} Three points form a triangle only if the sum of any two of the distances is greater than the third. Using the values above ($PQ\approx7.07$, $QR\approx7.21$, $PR\approx1.41$):
\begin{align*}
PQ+PR&\approx 7.07+1.41=8.48>QR\ (\approx7.21)\\
QR+PR&\approx 7.21+1.41=8.62>PQ\ (\approx7.07)\\
PQ+QR&\approx 7.07+7.21=14.28>PR\ (\approx1.41)
\end{align*}
All three hold, so $P,Q,R$ are not in a straight line and form a triangle.

\textbf{Step 3: what type of triangle?} By the converse of the Pythagoras theorem (if $a^2+b^2=c^2$ for the sides of a triangle, the angle opposite $c$ is $90^\circ$). The square of each side is the number under its square root in Step 1: $PQ^2=50$, $PR^2=2$, $QR^2=52$. The longest side is $QR$, so test it:
\[
PQ^2+PR^2=50+2=52=QR^2 .
\]
So the angle opposite $QR$, which is $\angle P$, is $90^\circ$. The sides $PQ,QR,PR$ are all different, so the triangle is not isosceles; it is a right-angled triangle.

\ans{The points form a triangle. It is a right-angled triangle, with the right angle at $P$.}

%---------------------------------------------------------------
\subsection*{Example 2}
\textbf{Question.} Show that the points $(1,7)$, $(4,2)$, $(-1,-1)$ and $(-4,4)$ are the vertices of a square.

\textbf{Solution.} Let $A(1,7)$, $B(4,2)$, $C(-1,-1)$, $D(-4,4)$.

A quadrilateral is a square if all four sides are equal and both diagonals are equal. By the distance formula $d=\sqrt{(x_2-x_1)^2+(y_2-y_1)^2}$, find the four sides and the two diagonals.

\distnote
\begin{align*}
\distn{AB}{1}{7}{4}{2}\\
\distn{BC}{4}{2}{-1}{-1}\\
\distn{CD}{-1}{-1}{-4}{4}\\
\distn{DA}{-4}{4}{1}{7}\\
\dist{AC}{1}{7}{-1}{-1}{2\sqrt{17}}\\
\dist{BD}{4}{2}{-4}{4}{2\sqrt{17}}
\end{align*}
Here $\sqrt{34}$ cannot be simplified, and $\sqrt{68}=\sqrt{4\times17}=2\sqrt{17}$.

Since $AB=BC=CD=DA=\sqrt{34}$ and $AC=BD=2\sqrt{17}$, all four sides are equal and the two diagonals are equal.

\ans{$ABCD$ is a square.}

\textbf{Alternative solution.} Using the converse of the Pythagoras theorem in $\triangle ADC$: $AD^2+DC^2=34+34=68=AC^2$ (the values $34$, $34$, $68$ are the numbers under the square roots above), so $\angle D=90^\circ$. A quadrilateral with all four sides equal and one angle $90^\circ$ is a square.

%---------------------------------------------------------------
\subsection*{Example 3}
\textbf{Question.} Fig.\ 7.6 shows the arrangement of desks in a classroom. Ashima, Bharti and Camella are seated at $A(3,1)$, $B(6,4)$ and $C(8,6)$ respectively. Do you think they are seated in a line? Give reasons for your answer.

\begin{center}
\ncertgrid{\fill (3,1) circle (3pt) node[above left]{$A$};
\fill (6,4) circle (3pt) node[above left]{$B$};
\fill (8,6) circle (3pt) node[above left]{$C$};}\\
{\small Fig.\ 7.6 (redrawn)}
\end{center}

\textbf{Solution.} The coordinates of the seats are given in the question: $A(3,1)$, $B(6,4)$, $C(8,6)$.

Three points are collinear if the largest of the three distances equals the sum of the other two. By the distance formula $d=\sqrt{(x_2-x_1)^2+(y_2-y_1)^2}$:

\distnote
\begin{align*}
\dist{AB}{3}{1}{6}{4}{3\sqrt2}\\
\dist{BC}{6}{4}{8}{6}{2\sqrt2}\\
\dist{AC}{3}{1}{8}{6}{5\sqrt2}
\end{align*}
Here $\sqrt{18}=\sqrt{9\times2}=3\sqrt2$, $\sqrt8=\sqrt{4\times2}=2\sqrt2$ and $\sqrt{50}=\sqrt{25\times2}=5\sqrt2$.

The largest distance is $AC=5\sqrt2$. Add the other two:
\[
AB+BC=3\sqrt2+2\sqrt2=(3+2)\sqrt2=5\sqrt2=AC .
\]
\ans{Since $AB+BC=AC$, the points $A$, $B$, $C$ are collinear. Ashima, Bharti and Camella are seated in a line.}

%---------------------------------------------------------------
\subsection*{Example 4}
\textbf{Question.} Find a relation between $x$ and $y$ such that the point $(x,y)$ is equidistant from the points $(7,1)$ and $(3,5)$.

\textbf{Solution.} Let $P(x,y)$ be equidistant from $A(7,1)$ and $B(3,5)$, so $AP=BP$ and therefore $AP^2=BP^2$ (squaring removes the square roots).

By the distance formula, $d^2=(x_2-x_1)^2+(y_2-y_1)^2$, so
\[
AP^2=(x-7)^2+(y-1)^2,\qquad BP^2=(x-3)^2+(y-5)^2 .
\]
Equate them and expand each bracket with $(a-b)^2=a^2-2ab+b^2$:
\begin{align*}
(x-7)^2+(y-1)^2&=(x-3)^2+(y-5)^2\\
(x^2-14x+49)+(y^2-2y+1)&=(x^2-6x+9)+(y^2-10y+25)\\
x^2+y^2-14x-2y+50&=x^2+y^2-6x-10y+34
\end{align*}
Cancel $x^2+y^2$ from both sides and collect the terms on the left:
\begin{align*}
-14x-2y+50&=-6x-10y+34\\
-14x+6x-2y+10y+50-34&=0\\
-8x+8y+16&=0\\
x-y&=2\qquad(\text{dividing by }-8)
\end{align*}
\slip{When expanding $(x-7)^2$ the middle term $-14x$ is often forgotten; $(x-7)^2\neq x^2-49$.}

\textbf{Verification.} The point $(4,2)$ satisfies $x-y=4-2=2$. Its distances from $A(7,1)$ and $B(3,5)$ are
\begin{align*}
\distn{AP}{7}{1}{4}{2}\\
\distn{BP}{3}{5}{4}{2}
\end{align*}
Both are $\sqrt{10}$, so the point is equidistant, as required.

\ans{The required relation is $x-y=2$.}

%---------------------------------------------------------------
\subsection*{Example 5}
\textbf{Question.} Find a point on the $y$-axis which is equidistant from the points $A(6,5)$ and $B(-4,3)$.

\textbf{Solution.} A point on the $y$-axis has the form $(0,y)$ (its $x$-coordinate is $0$). Let $P(0,y)$ be equidistant from $A$ and $B$, so $PA=PB$ and $PA^2=PB^2$.

By the distance formula $d^2=(x_2-x_1)^2+(y_2-y_1)^2$:
\[
PA^2=(6-0)^2+(5-y)^2,\qquad PB^2=(-4-0)^2+(3-y)^2 .
\]
Equate them and expand with $(a-b)^2=a^2-2ab+b^2$:
\begin{align*}
(6-0)^2+(5-y)^2&=(-4-0)^2+(3-y)^2\\
36+(25-10y+y^2)&=16+(9-6y+y^2)\\
61-10y+y^2&=25-6y+y^2\\
61-10y&=25-6y\\
36&=4y\\
y&=9
\end{align*}
So $P=(0,9)$.

\textbf{Verification.} With $P(0,9)$:
\begin{align*}
\dist{PA}{0}{9}{6}{5}{2\sqrt{13}}\\
\dist{PB}{0}{9}{-4}{3}{2\sqrt{13}}
\end{align*}
Both distances are $\sqrt{52}=2\sqrt{13}$.

\ans{The required point is $(0,9)$.}
